\section{Causal quantization with feedback and stopping}\label{sec:emulation}
\subsection{Admissible state domains}
Consider a predictive realization $s_t(h_t)$ with a known initial state
$s_0$. Let $(D_t,d_t)$ be admissible metric state domains containing every
attainable state. They may also contain normalized beliefs that are not
attained at that precise time. This causes no problem if all updates and
readouts are genuinely defined on $D_t$. What is forbidden is evaluating an
update at an artificial point where it has no experimental meaning or no
stated extension.

Assume common action availability on each domain, or work within a fixed
continuation-interface component. For each input $(a,y)$, let
$F_t(\cdot,a,y):D_t\to D_{t+1}$ be Borel and satisfy
\begin{equation}\label{eq:update-lip}
 d_{t+1}(F_t(s,a,y),F_t(\tilde s,a,y))\le L_td_t(s,\tilde s).
\end{equation}
The known-calibration formulas implementing $F_t$ are independent of the
unknown parameter. If there is a parameterized law, write the next-report
kernel as $g_t^\lambda(dy\mid s,a)$ and assume
\begin{equation}\label{eq:report-lip}
 \norm{g_t^\lambda(\cdot\mid s,a)-g_t^\lambda(\cdot\mid\tilde s,a)}_{\TV}
 \le\kappa_td_t(s,\tilde s)
\end{equation}
uniformly over the relevant parameter domain. The same Borel state update
must represent every law for a parameter-uniform statement. For a single
prepared Bayesian experiment there is no additional parameter supremum.

Choose finite representative sets $C_t\subset D_t$ with $|C_t|\le M_t$
and Borel quantizers $q_t:D_t\to C_t$ satisfying
$d_t(s,q_ts)\le r_t$. A numerical update
$\widetilde F_t(c,a,y)\in D_{t+1}$ may replace the exact update at a
representative, with
\begin{equation}\label{eq:update-numerics}
 d_{t+1}(\widetilde F_t(c,a,y),F_t(c,a,y))\le\eta_{t+1}.
\end{equation}
Set $\hat s_0=s_0$ and implement
\begin{equation}\label{eq:finite-machine}
 \hat s_{t+1}=q_{t+1}\widetilde F_t(\hat s_t,a_t,y_{t+1}).
\end{equation}
The persistent state is only the index of $\hat s_t$ in $C_t$.

\begin{theorem}[Finite-state emulation]\label{thm:emulation}
Under \eqref{eq:update-lip}--\eqref{eq:update-numerics}, define
\begin{equation}\label{eq:error-b}
 b_0=0,\quad b_{t+1}=L_tb_t+r_{t+1}+\eta_{t+1},\quad
 b_t=\sum_{j=1}^t(r_j+\eta_j)\prod_{i=j}^{t-1}L_i.
\end{equation}
Then the following assertions hold.

\textup{(i)} On every admitted common command/report sequence,
$d_t(s_t,\hat s_t)\le b_t$. If the physical prediction readout is
$K_t$-Lipschitz, its excess squared error is at most $K_t^2b_t^2$.

\textup{(ii)} Define the finite-state reference experiment to use the same
representative update \eqref{eq:finite-machine}, but to draw its next
report from $g_t^\lambda(\cdot\mid\hat s_t,a_t)$. For any common causal
feedback controller, the visible path laws of the actual experiment with
the representative machine and of this finite-state reference experiment
satisfy
\begin{equation}\label{eq:emulation-TV}
 \norm{P_{N,\lambda}^{\alpha}-\widehat P_{N,\lambda}^{\alpha}}_{\TV}
 \le \min\left\{1,\sum_{t=0}^{N-1}\kappa_tb_t\right\}.
\end{equation}
The histories here include the applied actions, reports, and retained
labels. A shared visible stopping rule replaces the deterministic sum by
$\E_{P^\alpha_\lambda}\sum_{t<\sigma\wedge N}\kappa_tb_t$.

\textup{(iii)} Any additional causal simulator is charged as in
Theorem~\ref{thm:composition}. A feedback controller with $D_t$ own states
is implemented with at most $D_tM_t$ persistent pairs. For a common bounded
path-dependent loss, the risk-transfer error is at most its range times the
right side of \eqref{eq:emulation-TV}.
\end{theorem}
\begin{proof}
For \textup{(i)}, insert successively the exact update from the retained
representative and its numerical approximation. The triangle inequality,
\eqref{eq:update-lip}, the numerical error, and the covering radius give
\[
 d_{t+1}(s_{t+1},\hat s_{t+1})
 \le L_td_t(s_t,\hat s_t)+\eta_{t+1}+r_{t+1}.
\]
Induction proves \eqref{eq:error-b}. The physical readout bound follows
by its Lipschitz property and the conditional squared-loss identity.
This argument remains valid when the actions depend on the retained label:
it compares the exact posterior and the representative along the actual
same sequence of applied commands and reports.

For \textup{(ii)}, at any common visible prefix the applied controller has
the same action kernel. The representative $\hat s_t$ is the same
measurable function of that prefix in both experiments. In the actual law,
the exact predictive state is $s_t$; in the finite-state reference law, the
next-report kernel is evaluated at $\hat s_t$. Their conditional
variation distance is at most $\kappa_tb_t$ by \textup{(i)} and
\eqref{eq:report-lip}. The next label is a common deterministic function
of the action, report, and current label, so appending it does not increase
the conditional variation distance. Apply Lemma~\ref{lem:telescoping}.
Stopping makes the conditional kernels identical after the stopping time,
which gives the predictable active-step bound.

For \textup{(iii)}, pair the retained representative with the controller's
internal state. Each next action and update then reads only this pair and
current inputs. The state cardinality is bounded by the product. The
bounded-loss and extra-simulator statements are precisely the conclusions
of Theorem~\ref{thm:composition}.
\end{proof}

The theorem is a directional emulation statement. It does not construct a
reverse kernel from the finite-state reference experiment to the raw
physical data. It therefore does not claim two-sided Le Cam equivalence.
Nor does it estimate an unknown predictive update. The implemented update
and representatives are known; uniformity over a statistical family is
available only under the explicitly common update conditions above.

\subsection{What feedback does and does not require}
There are two different uses of a finite filter. One can predict on an
arbitrarily selected actual input sequence; this uses only
Theorem~\ref{thm:emulation}\textup{(i)}. One can also use a compressed
controller to choose inputs. That still gives accurate filtering along its
own actual run, but comparison with the original full-information optimal
controller is another problem. Part \textup{(ii)} compares a common
controller across two experiments and so supplies a valid route through
experiment simulation. It does not require the controller to be Lipschitz,
because identical prefixes produce identical action laws. A discontinuous
policy evaluated at two different approximate states is not being assumed
to choose the same action.

For example, minimizing a terminal common decision loss over all
$D_t$-state controllers in the finite-state reference experiment gives a
source controller with $D_tM_t$ states and the corresponding
\eqref{eq:emulation-TV} risk allowance. This need not approximate the
unrestricted optimal source policy, whose comparison would require a
separate richness theorem for the controller class.

\subsection{Uniform prediction and finite-horizon path comparison}
\begin{corollary}[Contractive implementation]\label{cor:contractive}
If $L_t\le\rho<1$, $r_t\le r$, and $\eta_t\le\eta$ uniformly, then
\[
 \sup_td_t(s_t,\hat s_t)\le\frac{r+\eta}{1-\rho}.
\]
For a bounded physical readout with uniform Lipschitz constant $K$, the
prediction excess is at most $K^2(r+\eta)^2/(1-\rho)^2$ at every time.
If also $\kappa_t\le\kappa$, the path-law error on $N$ steps is at most
$\min\{1,N\kappa(r+\eta)/(1-\rho)\}$.
\end{corollary}
\begin{proof}
Sum the geometric series in \eqref{eq:error-b}, then apply the two parts of
Theorem~\ref{thm:emulation}.
\end{proof}

A uniform per-checkpoint prediction bound is not uniform total-variation
closeness of two infinite path laws. The last display may grow with $N$.
Two processes that are arbitrarily close at each step can remain
statistically distinguishable from a sufficiently long record.
If a general construction uses a different arbitrary codebook at each
time, it is a time-indexed machine description, not automatically a finite
uniform program. The dynamic construction in the next section uses one
fixed invariant domain and one fixed codebook, so that distinction does
not obstruct its infinite-time statement.

\subsection{Model or calibration perturbations}
Suppose an implementable reference update $F_t^0$ has
$d(F_t(s,a,y),F_t^0(s,a,y))\le\delta_{t+1}$ on the relevant domain, and its
report law differs by at most $\zeta_t$ at the same state. Repeating the
proof replaces $r_j+\eta_j$ by $r_j+\eta_j+\delta_j$ in
\eqref{eq:error-b} and $\kappa_tb_t$ by $\kappa_tb_t+\zeta_t$ in
\eqref{eq:emulation-TV}. This is a deterministic certificate, not an
unpaid pilot. If a preliminary estimate yields such bounds only on an event
of probability $1-p$, bounded risks receive an additional allowance of at
most their range times $p$, and the pilot's data and persistent state must
be charged.

The perturbation must remain in the same physical metric and admissible
domain. A reparametrization with an unbounded inverse near a degeneracy
cannot be inserted while keeping the old constants.
